\subsection{The nonconstant six-state four-row branch}

\begin{lemma}[Six bounded rank-two action hulls]
\label{C:six-rank-two-bounded-hulls}
Consider the six rank-two $4\times4$ action hulls numbered
\[
 137,169,172,188,222,225
\]
in Proposition~\ref{B:four-by-four-product25}.
Delete every actor on whose eventual cyclic group one of the
two stored integral degree-two classes restricts nontrivially.
Among closed remaining submonoids containing the band and identity,
of order at most twenty-two and rational degree-two rank two,
the following is a complete independently checked count.
\[
\begin{array}{c|r|r|r}
\text{hull}&\text{closed subsets}&\text{orbits}&
 (b_2,b_3,b_4)_{\mathbb F_2}=(2,0,\text{positive})\\\hline
137&332&168&4\\
169&570&157&0\\
172&708&367&0\\
188&332&168&4\\
222&7&4&0\\
225&7&4&0
\end{array}
\]
The third column uses automorphisms and anti-isomorphisms.
All eight representatives in the last column have order twenty.
In either orientation, if their outside actors preserve a two-element
subset of the four column labels, then the identity profile is the
only outside profile whose restriction to that subset is nonconstant.
\end{lemma}
\begin{proof}
The independent fixed-point closure program
\path{code/a_rank_two_act_closure.cpp} starts with the sixteen constant
profiles and the identity, adjoins one allowed actor at a time, and
repeatedly adjoins every product until closure.  It rejects a closure
only when it contains a forbidden actor or more than five optional
actors.  Thus every allowed closed subset in the stated size range is
reached.  Rational elimination on the nine balanced-tensor generators
retains exactly relation rank seven.  It uses primitive integer rows;
the original row squared norms are at most eight, so the nine-column
minor bound also bounds the intermediate integer arithmetic.

The separate checker \path{code/a_rank_two_act_audit.py} reconstructs
every literal actor product, every associativity equation, and every
balancing relation.  It independently verifies the cyclic-character
exclusions, the input allowed-actor bits, the complete closure sets,
and the disjoint orbit coverage under every supplied automorphism or
anti-isomorphism.  It then computes the three contracted two-sided
bar differentials over $\mathbb F_2$ in an independent tuple ordering,
checking all $827{,}910$ relevant $d^2=0$ equations.  The band double
shift identifies their homology with degrees two, three and four.
All $868$ representative results agree with the original producer.
The exact record is
\path{code/outputs/a_rank_two_act_audit.json}, whose completed flag is
true; the six individual closure logs are retained alongside it.

Finally the checker tests all six two-column subsets in both
orientations of every surviving representative.  With the original
zero-based case labels, the only possibilities are
\[
\begin{array}{c|c|c|c}
\text{hull}&\text{case}&\text{orientation}&\text{old columns}\\\hline
137&152,169&\text{ordinary}&\{2,3\}\\
188&168,204&\text{opposite}&\{0,2\}.
\end{array}
\]
The other cases, $279,291$ in hull $137$ and $208,210$ in hull $188$,
have no such subset in either orientation.  In each displayed case
the sole nonconstant outside profile on the subset is the identity,
whose row map is also the identity.  Testing both orientations is
necessary because the orbit reduction permits anti-isomorphisms.
\end{proof}

\begin{theorem}[Nonconstant six-state four-row acts]
\label{C:four-row-nonconstant-six-states}
Let $M$ have minimal band $4\times2$ and at most six outside elements.
If a finite right act has at most six states and some minimal-band
actor is nonconstant, its act cone does not have type
$\mathbb CP^2$.
In particular, the nonconstant case of the twenty-one-arrow profile
$(|M|,|S|)=(14,6)$ is excluded.
\end{theorem}
\begin{proof}
Suppose the cone $X$ is a projective plane.  A state fixed by the
minimal band would make all its degree-two products zero, so there
is no such state.  The act is connected.  Every core orbit therefore
has two points and is individually $M$-invariant.  Nonconstant band
action requires at least two core orbits, and they cannot exhaust
the connected act.  At most four states are impossible.

For five states there are two core orbits.  The projective-orbit
kernel and finite-quotient model lemmas give a compressed monoid
with minimal band $4\times3$, a finite four-dimensional model,
$H_2=\mathbb Z^2$, $H_4=\mathbb Z$, and no other positive homology.
The rank-two conclusion in the proof of
Theorem~\ref{B:three-by-four-targets}, after opposition, kills all
its rational degree-two products.  Its degree-two classes pull back
onto those of $X$, a contradiction.

For six states there are again exactly two core orbits: three would
exhaust the act and disconnect it.  The constant-action quotient
has four states.  Its compression $N$ has a $4\times4$ minimal band,
the same $q\leq6$ outside elements, and therefore order at most
twenty-two.  Lemmas~\ref{C:two-column-projective-orbit-kernel} and
\ref{C:two-column-finite-quotient-model} give
\[
 H_2(BN;\mathbb Z)=\mathbb Z^2,\qquad
 H_4(BN;\mathbb Z)=\mathbb Z,
\]
zero other positive homology, and an actual finite four-dimensional
CW model.  Moreover $H^2(BN;\mathbb Q)\to H^2(X;\mathbb Q)$ is
surjective.  Thus a square on $BN$ must have nonzero pullback.

Use the complete rank-two hull classification in the proof of
Proposition~\ref{B:four-by-four-product25}.  The rank-one-form,
repeated-point, zero-fourth-homology and exact-zero-product cases
have all rational degree-two products zero, and so cannot occur.
The three common-radical hulls numbered $161,237,400$ also have all
products zero.  The separately verified orbit-$346$ subset argument
gives zero fourth rational homology for every possible image of
order at most twenty-four.  Consequently only the six hulls in
Lemma~\ref{C:six-rank-two-bounded-hulls} remain.  We do not exclude
them merely from the existence of a radical line: the map to $X$
could kill that line.

Let $H$ be the faithful row-and-column image of $N$ in one of those
hulls.  This quotient fixes its band pointwise.  It preserves
integral $H_2$, by the identical balanced-tensor presentation, and
is onto on $H_3(-;\mathbb F_2)$ by
Lemma~\ref{B:band-quotient-H3-surjection}.
Thus $b_2(H;\mathbb F_2)=2$ and $b_3(H;\mathbb F_2)=0$.
If $b_4(H;\mathbb F_2)=0$, its rational fourth homology is zero,
and the isomorphism on rational $H^2$ again kills the source products.
The forbidden cyclic-character actors cannot occur: a local cyclic
group in the image lifts surjectively from an eventual local cyclic
group of $N$, and a nonzero integral character remains nonzero on
that lift.  Miller's theorem makes every such restriction zero for
the finite-dimensional source $BN$.
Lemma~\ref{C:six-rank-two-bounded-hulls} therefore applies to $H$.

Let $T$ be the image in the four quotient states of the original two
minimal columns.  If $|T|=1$, every old column term $b_m$ in the
transferred affine cocycles is zero, which already kills all products.
Otherwise $T$ is a two-element subset preserved by every outside
actor.  This invariance passes literally to the profile image $H$.
The lemma says that an image profile nonconstant on $T$ must be the
identity profile.  Any original outside actor mapping to that
profile has the identity row map; an actor mapping to a band profile
has constant columns.  Hence every actor of the original $M$ with
nonconstant old column map fixes every row.
The common-fixed-row cup primitive now makes all products on $X$
zero, the final contradiction.

This last argument explicitly permits duplicate outside profiles.
No conclusion about the higher homology or finite type of a proper
image is substituted for cup-product naturality.
\end{proof}
